\section{Limitations and Conclusion}
\label{sec:main-limitations}

\begin{itemize}
\item The theory is about explicit information structures or restricted hard-route classes.
\item Full backpropagation through the representation is outside the repair-only lower bound.
\item Soft-DGM requires candidate recall and margin assumptions to approximate the hard route.
\item Held-out guarantees require fresh or reusable validation data rather than adaptive reuse of the same scoring buffer.
\item The experiments are diagnostics plus pure-DGM pilots, not broad SOTA evidence.
\end{itemize}
The non-oracle Image-HQS result is especially narrow: it supports local posterior repair under learned visual concepts, but it does not establish exact hierarchy recovery from images or a complete memory-latency frontier.
The pure-DGM pilots show that no-backprop memory is viable when the fixed geometry is adequate, and the CIFAR-10 handcraft runs show that Pure DGM can improve sharply once the fixed distinction basis has basic visual invariances.
They also show the current limit: raw pixels and PCA are not enough, naive ``store more views'' scaling saturates without better memory selection, and the current FullDGM-refine implementation still trails a ridge readout on the same fixed HOG/color/DCT/LBP memory views.
The trainable-CNN and character-LM pilots are therefore treated as non-pure references rather than the main claim.
The next empirical step is a budget-matched comparison against stronger tree, RLS, $k$NN, handcrafted-feature SVM/ridge, cache, adapter, and prototype-growth baselines under the same prequential contract.
The method also needs better policies for merging, eviction, nonlinear distinctions, and adaptive validation if it is to be used as a long-running system rather than a diagnostic implementation.

The main point is the separation between two runtime operations.
Repair improves predictions inside the current quotient; refinement changes the quotient when consequences reveal a missing distinction.
DGM gives one guarded graph-memory implementation of that decision: accept consequence-valued distinctions as edges, otherwise repair local memories.
That narrower framing is the intended contribution: a decision-theoretic obstruction, a value criterion for distinctions, and a guarded graph-memory mechanism aligned with those claims.
